% Physical unit 007; accepted Session 02.
\GordanUnitStart{007}{printed page 2}
\GordanTopFolio{2}
of this circle makes with the equator an angle equal to all the others; this angle is called the \gq{reduced latitude} of the parallel circle of the spheroid and is denoted by the sign \gq{$l$}. ---

\begin{center}\S\ 6.\end{center}
Thus every circle of the spheroid, since it is determined both by its latitude \gq{$B$} and by the latitude \gq{$l$}, is called the circle \gq{$B$}, or the circle \gq{$l$}. The radius of the circle \gq{$l$} is $a\cos l$.

\begin{center}\S\ 7.\end{center}
Either latitude is easily determined from the other. In fact, between the angle $l$ and the angle $B$ there is the following relation:
\[
  a\,\tg l=b\,\tg B.
\]
Or, since the eccentricity of the meridian ellipse, namely $\sqrt{1-\frac{b^2}{a^2}}$, is denoted by \gq{$e$}, the relation is:
\[
  \tg l=\sqrt{1-e^2}\,\tg B.
\]
Although this is the simplest way of reducing the latitudes $B$ and $l$ one to the other, we shall nevertheless often use other relations, all of which are known to be reducible to the preceding one.

The most important of these relations are:
\[
\begin{aligned}
\sin l&={b\over a}{\sin B\over\sqrt{1-e^2\sin^2 B}},
&\qquad
\sin B&={\sin l\over\sqrt{1-e^2\cos^2 l}},\\
\cos l&={\cos B\over\sqrt{1-e^2\sin^2 B}},
&
\cos B&={b\over a}{\cos l\over\sqrt{1-e^2\cos^2 l}},\\
{b^2\over a^2}&=1-e^2=(1-e^2\sin^2 B)(1-e^2\cos^2 l).
\end{aligned}
\]
Finally, putting
\[
  {a-b\over a+b}={1-\sqrt{1-e^2}\over 1+\sqrt{1-e^2}}=\varepsilon,
\]
we have
\[
  B=l+\varepsilon\sin 2l+{\varepsilon^2\over 2}\sin 4l+{\varepsilon^3\over 3}\sin 6l.
\]

\begin{center}\S\ 8.\end{center}
Since, as we said above, every point of the spheroid is determined by the angles $l$ and $\psi$, these may be taken as coordinates. Just as the position of any point in a plane is defined by two coordinates, and a plane curve is determined by a relation between those two coordinates, so a relation between $l$ and $\psi$ expresses a curve on the spheroid. Further, as in the discussion of plane curves one usually treats most especially the length and the inclination at each point to some fixed axis
